\documentclass[10pt,a4paper]{amsart}
\usepackage[margin=19mm]{geometry}
\usepackage{amsmath,amssymb,mathtools}
\usepackage[scr=rsfs,cal=euler]{mathalfa}
\usepackage{xfrac}
\usepackage[unicode,hidelinks,bookmarks=false]{hyperref}
\newcommand{\ie}{i.e.\ }
\newcommand{\cf}{cf.\ }
\newcommand{\resp}{respectively\ }
\newcommand{\Subsection}{\subsection}
\providecommand{\nobreakdash}{\nobreak\hbox{-}}
\setlength{\emergencystretch}{2em}
\multlinegap0pt
\newtheorem{lemma}{Lemma}
\newcommand{\Zpos}[1]{{#1}_{\zar}^{\leq}}
\DeclareMathOperator{\zar}{zar}
\title{\Large On Galois categories and condensed contractible schemes}
\author{Catrin Mair}
\date{Source: arXiv:2605.10358v1 (CC BY 4.0)}
\begin{document}
\maketitle

\noindent\emph{Source and licence.} \Large On Galois categories and condensed contractible schemes. Version arXiv:2605.10358v1 (CC BY 4.0), distributed by arXiv under the Creative Commons Attribution 4.0 International licence, http://creativecommons.org/licenses/by/4.0/, as recorded in the arXiv licence metadata for this exact version and shown on the abstract page https://arxiv.org/abs/2605.10358v1. Full licence notice: \texttt{licenses/arxiv-2605.10358-CC-BY-4.0.txt}.\par\smallskip

\noindent\textit{Editorial scope.} Counted unit: the article's lemma lemma (source0.tex lines 1655--1662). The article's complete original proof of it is delivered with it (source0.tex lines 1663--1680, one \texttt{proof} environment of 18 lines). No mathematics is written, repaired or re-derived: the statement and the proof are the article's own lines, in the article's own notation, and no retained line is substituted or re-flowed.  No further source-local context is needed: every cross-reference of every retained line resolves inside the two delivered spans.  Nothing was omitted from any retained span, so the package carries no omission record.  No retained line contains a citation, so the wrapper carries no bibliography and no bibliography entry.  No retained line contains a figure environment, an image-inclusion command or drawing code, so no image is delivered and the package declares no asset dependency.  Editorial formatting only, in addition: an AMS \texttt{amsart} preamble that restates the article's own declarations and macros used by the retained lines, namely the environments \texttt{lemma} on separate counters, together with the standard packages those lines need.  The retained environments print in this document as Lemma 1 in order of appearance, whereas the article prints the same items with the numbering of its own full document.  The complete native source of the article as distributed in the arXiv e-print for this exact version is bundled unchanged as \texttt{sources/200291/source0.tex}.
\begin{lemma}\label[lemma]{lem:XZarconditions} Let $X$ be a qcqs scheme.
The following assertions are equivalent.
\begin{itemize}
\item[a)] The Zariski poset $ \Zpos{X}$ has a maximal element (respectively, a minimal element).
\item[b)] The Zariski poset $ \Zpos{X}$ is (co)directed.
\item[c)] The scheme $X$ is irreducible (respectively, local, i.e., corresponds to a local ring).
\end{itemize} 
\end{lemma}

\begin{proof}
The equivalences of a), b) and c) rely on the definition of the order by specializations.
For a)$\implies b)$ note that a poset with a maximal or minimal object is, in particular, (co)directed. For  b) $\implies$ c), we have to argue why there is a unique generic (closed) point:
The poset $ \Zpos{X}$ to be (co)directed means that any two points are specializations (generalizations) of a common point.
We handle the two cases separately.
\begin{itemize}
    \item[\textbf{1.}] \textbf{The directed case:}
    To satisfy the condition on specializations, the sober topological space $|X|$ cannot have more than one irreducible component since two maximal points, i.e., the unique generic points of irreducible components, do not specialize further.
    This implies irreducibility.
    \item[\textbf{2.}] \textbf{The codirected case:}
    Every closed point of $|X|$ does not generalize further as there is no other point lying in its closure. Hence, a scheme with more than one closed point contradicts the condition on $ \Zpos{X}$ to be codirected.
    A scheme that is quasi-compact and contains exactly one closed point is known to be equivalent to the spectrum of a local ring.
    This follows from the fact that every quasi-compact scheme, and thus also every closed subset of a quasi-compact scheme, has a closed point.
    Then the complement of an affine open containing the unique closed point has to be empty.
    This already implies that the scheme itself is affine with a unique closed point corresponding to a unique maximal ideal of the underlying ring.
    \end{itemize}
From c) to a), it is clear that the unique generic point or, respectively, the unique closed point of the scheme constitutes a maximal or, respectively, a minimal element of $ \Zpos{X}$.
\end{proof}

\end{document}
